\documentclass[11pt]{article}

\usepackage[a4paper,margin=1in]{geometry}
\usepackage{amsmath,amssymb,amsthm,mathtools}
\usepackage{enumitem}
\usepackage{hyperref}

\title{Research Plan:\
From Lucas Digit Geometry to Generalized Carry Calculus}
\author{}
\date{\today}

\newtheorem{question}{Research Question}
\newtheorem{expectation}{Expected Outcome}

\begin{document}

\maketitle

\begin{abstract}
This document records the current research plan following the development of a digit-structured theory of Lucas HIT sets and a corresponding local digit calculus. The central objective of the next phase is to determine which parts of the construction form general mathematical structures independent of the original parameterization, and separately to investigate whether any factor-sensitive information can be extracted when the prime radix is unknown. The research is divided into a mathematical generalization track and a high-risk factorization track.
\end{abstract}

\section{Current State}

Let

$$
m=\sum_{i\geq 0}m_i p^i
$$

be the base-$p$ expansion of $m$. Lucas' theorem gives

$$
p\nmid \binom{m}{t}
\iff
t_i\leq m_i\quad\text{for every }i.
$$

Thus the MISS set is the digitwise box

$$
\operatorname{MISS}_p(m)
=
\{t:0\leq t\leq m,\ t_i\leq m_i\},
$$

and the HIT set is its complement in $[0,m]$.

The number of MISS values is

$$
M(m)=\prod_i(m_i+1),
$$

and therefore

$$
H(m)
=
|\operatorname{HIT}_p(m)|
=
m+1-\prod_i(m_i+1).
$$

The ordering of the MISS set induces a mixed-radix rank

$$
R(t)=\sum_i t_iW_i,
\qquad
W_i=\prod_{j<i}(m_j+1).
$$

Consecutive MISS values produce HIT intervals whose carry structure is governed by

$$
G_r
=
p^r-1-\sum_{i<r}m_ip^i,
$$

with multiplicity

$$
C_r
=
m_r\prod_{i>r}(m_i+1).
$$

This structure has subsequently been converted into a local digit calculus with exact prefix evaluation and path-independent digit updates.

The computational work has also shown that several apparently promising factorization constructions were either ordinary multiplicative-order effects, radix artifacts, or implementation artifacts. Consequently, future experiments will be required to satisfy a stricter factor-blindness criterion.

\section{Research Track I: General Lucas Digit Geometry}

\subsection{Arbitrary $m$}

The first objective is to remove the special parameterization

$$
m=s_0+qp^e-1
$$

from the central theory.

\begin{question}
Does the complete interval decomposition, mixed-radix rank, carry-gap classification, and prefix calculus hold for every integer $m$ directly from its base-$p$ digits?
\end{question}

The expected general structure is

$$
\operatorname{HIT}_p(m)
=
\bigsqcup_{r,j} I_{r,j},
$$

where the intervals are generated by successive states of the mixed-radix MISS lattice.

\begin{expectation}
The special parameterized family will become a structured corollary of a more general arbitrary-digit theorem.
\end{expectation}

The main expected result is therefore not another special closed form, but a theorem explaining the entire one-dimensional interval grammar for arbitrary digit vectors.

\subsection{Generating Functions and Moments}

The MISS set has the generating function

$$
\sum_{t\in\operatorname{MISS}_p(m)}z^t
=
\prod_i
\left(
1+z^{p^i}
+\cdots+
z^{m_ip^i}
\right).
$$

Consequently,

$$
\sum_{t\in\operatorname{HIT}_p(m)}z^t
=
\frac{1-z^{m+1}}{1-z}
-
\prod_i
\frac{1-z^{(m_i+1)p^i}}
     {1-z^{p^i}},
$$

with the value at $z=1$ understood by continuity.

This gives a systematic route to

$$
\sum_{t\in\operatorname{HIT}}t^k
$$

and potentially to Fourier or other spectral transforms.

\begin{expectation}
The digit-product structure will yield exact closed expressions for HIT moments and transforms, extending the cardinality identity.
\end{expectation}

\subsection{Kummer Carry Geometry}

Lucas detects whether at least one carry occurs. Kummer gives the exact valuation:

$$
v_p\binom{m}{t}
=
\text{number of carries in }t+(m-t).
$$

The next goal is to replace the Boolean HIT/MISS classification by a carry-count state system:

$$
\mathcal{K}_k(m)
=
\left\{
t:
v_p\binom{m}{t}\geq k
\right\}.
$$

Rather than assuming that $\mathcal{K}_k(m)$ is always a refinement into ordinary intervals, the investigation will construct an explicit carry-state automaton.

\begin{expectation}
The mixed-radix calculus will generalize from a Boolean digit condition to a finite-state carry calculus capable of counting and locating higher $p$-adic valuation layers.
\end{expectation}

\subsection{Multinomial Generalization}

For multinomial coefficients,

$$
\binom{m}{t_1,\ldots,t_d},
$$

the no-carry condition becomes

$$
t_{1,i}+\cdots+t_{d,i}\leq m_i
$$

at every digit.

The one-dimensional digit box is therefore replaced by a local lattice simplex.

The local number of admissible digit tuples is

$$
\binom{m_i+d-1}{d-1}.
$$

This suggests a higher-dimensional analogue of the MISS product:

$$
\#\operatorname{MISS}_{d}(m)
=
\prod_i
\binom{m_i+d-1}{d-1},
$$

for the corresponding multinomial parameter domain.

\begin{expectation}
The interval grammar will generalize into a higher-dimensional lattice or polyhedral decomposition, with local digit-simplex factors replacing $(m_i+1)$.
\end{expectation}

\section{Research Track II: Factor-Sensitive Experiments}

The mathematical calculus currently assumes that the true prime radix $p$ is known. A factoring application would require an observable computable without knowing either hidden factor of

$$
N=pq.
$$

This imposes the following rule:

\begin{center}
\textbf{Every proposed factoring observable must be factor-blind.}
\end{center}

That is, its construction may depend on $N$, a trial base $b$, and the digit calculus induced by $b$, but not on $p$, $q$, $p-1$, $q-1$, or multiplicative orders computed using the hidden factors.

\subsection{Trial-Base Sensitivity}

For a trial base $b$, define a digit-calculus observable

$$
F_b(N).
$$

The central question is whether the behavior of $F_b(N)$ changes when $b$ has an arithmetic relationship with $N$, particularly when

$$
b\mid N
$$

or, in the semiprime setting, when $b$ equals a true prime factor.

The intended tests include differences such as

$$
F_b(N)-F_c(N)
$$

and

$$
\gcd\bigl(F_b(N)-F_c(N),N\bigr).
$$

The critical comparison is between

$$
b\mid N
$$

and

$$
b\nmid N.
$$

\begin{expectation}
Either a factor-sensitive base transition will appear, providing a possible route toward factorization, or the relevant observables will remain algebraically insensitive to the factor structure.
\end{expectation}

This is currently the highest-priority factorization experiment because it tests a genuine hypothesis without assuming that path independence itself must break.

\subsection{Local-Slice Factor-Blind Transforms}

The exact digit slices

$$
\mathcal{S}_r(m)
$$

provide structured subsets without explicitly scanning the entire interval.

Future experiments will investigate aggregate transforms of these slices, such as sums, products, generating functions, and polynomial evaluations.

The requirement is that any candidate quantity be formed without inserting the unknown factor itself.

\begin{expectation}
If the local slices contain factor-sensitive information, it should appear as an arithmetic asymmetry in a factor-blind aggregate rather than merely as another representation of the known digit structure.
\end{expectation}

\section{Research Track III: Discrete Calculus}

The local digit calculus admits commuting updates

$$
\Delta_r\Delta_s
=
\Delta_s\Delta_r
$$

and can be interpreted as a discrete potential system on a mixed-radix lattice.

Further work will classify higher differences

$$
\Delta_{r_1}\cdots\Delta_{r_k}F
$$

and determine when they vanish identically.

The aim is to identify a hierarchy analogous to discrete polynomial degree.

\begin{expectation}
Higher-order differences will exhibit systematic cancellation laws determined by the digit-local algebra, yielding structural theorems even if they provide no factoring algorithm.
\end{expectation}

However, vanishing of a discrete derivative by itself will not be interpreted as factor-sensitive. A factoring claim requires an explicit asymmetry

$$
p\mid X,
\qquad
q\nmid X,
$$

with $X$ constructed without knowledge of $p$ and $q$.

\section{Research Methodology}

Several recent experiments have shown that apparently strong factorization results can arise from trivial mechanisms. Future experiments will therefore distinguish carefully between:

\begin{enumerate}[label=(\roman*)]
\item structural identities;
\item accidental arithmetic collisions;
\item factor-order or Fermat effects;
\item representation artifacts;
\item genuine factor-sensitive observables.
\end{enumerate}

In particular, an observed GCD

$$
\gcd(X,N)\in\{p,q\}
$$

will not by itself be considered evidence of a new mechanism. The construction of $X$ must first be shown to be independent of the hidden factors, and known explanations such as multiplicative orders, $p-1$ smoothness, and radix valuation effects must be excluded.

\section{Expected Research Outcome}

There are two deliberately separate success criteria.

\subsection*{Mathematical success}

The main mathematical objective is to establish a general framework

$$
\boxed{
\text{Lucas digit order}
\longrightarrow
\text{mixed-radix geometry}
\longrightarrow
\text{interval/slice calculus}
\longrightarrow
\text{Kummer and multinomial generalizations}.
}
$$

This direction remains valuable independently of factorization.

\subsection*{Algorithmic success}

The high-risk objective is to determine whether a factor-blind trial-base observable exists such that

$$
F_b(N)
$$

changes in a detectable way when the trial base interacts with a hidden factor of $N$.

A successful result would provide a bridge from

$$
\text{digit calculus}
$$

to

$$
\text{factor-sensitive arithmetic}.
$$

A persistent negative result would nevertheless be valuable if it establishes that the present calculus is representation-rich but factor-neutral under broad classes of observables.

\section{Immediate Experimental Order}

The planned order of investigation is:

$$
\boxed{
\begin{array}{cl}
1. & \text{Arbitrary-$m$ interval theorem},\\[2mm]
2. & \text{Trial-base sensitivity},\\[2mm]
3. & \text{Kummer carry-state calculus},\\[2mm]
4. & \text{Multinomial generalization},\\[2mm]
5. & \text{Generating functions and spectral transforms},\\[2mm]
6. & \text{Higher discrete-difference structure}.
\end{array}
}
$$

The first direction strengthens the existing theory. The second is the principal high-risk factoring experiment. The remaining directions expand the mathematical framework regardless of whether factorization succeeds.

\section{Closing Perspective}

The project has reached a useful transition point.

The central question is no longer whether Lucas' theorem produces complicated HIT patterns. It does, and those patterns can already be described exactly.

The next question is whether this exact digit geometry has consequences beyond representation:

$$
\boxed{
\text{What new mathematics follows from the mixed-radix structure?}
}
$$

and, separately,

$$
\boxed{
\text{Can any of that structure remain informative when the true radix is unknown?}
}
$$

The future research program is therefore designed to pursue both possibilities without conflating structural elegance with factorization evidence.

\end{document}
